% ===================================================================
%  TABELL Nr 10 — Havet. — Hamnen.
%  Traduction 2026 d'apres texto/io, controlee sur texto/fr et
%  texto/es. Consigne : voir texto/sv/10-tabell-01.tex.
%
%  Le francais decale sa numerotation d'un cran a partir de l'alinea
%  8 ; c'est une coquille de l'atelier francais, et la colonne suit
%  l'ido, qui compte juste.
% ===================================================================

%%K t10-apar-1 apar
\VUsaut{4.0mm}
\VUtitre{15.0pt}{60}{TABELL Nr 10}
\VUsaut{3.0mm}
\VUpk{11.4pt}{40}{Havet. --- Hamnen.}
\VUsaut{3.0mm}

%%K t10-01-1 p
1. --- Om sommaren, medan somliga söker lugn och svalka på ett berg, föredrar andra att fara
till kusten. Så gjorde vi förra året, för vi tycker mycket om
\VUgras{Oceanen} \textsuperscript{(1)}.

%%K t10-02-1 p
2. --- Vad mig beträffar, känner jag ett stort nöje av att betrakta den ofantliga vattenvidd
som sträcker sig ända till \VUgras{horisonten} \textsuperscript{(2)}, och av att se de enorma
\VUgras{ångfartygen} \textsuperscript{(3)} fara i väg på en lång \VUgras{överfart} med de
resande som far till Amerika ombord.

%%K t10-03-1 p
3. --- Därför väljer min far, som känner min böjelse, alltid en mycket lugn strand att
tillbringa ferierna på, som ligger vid en av våra stora
\VUgras{örlogshamnar}.

%%K t10-03-2 p
Under vår sista vistelse kom \VUgras{eskadern} och ankrade bakom den enorma
\VUgras{vågbrytaren} \textsuperscript{(4)} som stänger och skyddar
\VUgras{örlogshamnen} \textsuperscript{(5)}, och det lyckades mig som jag önskade
beundra de olika \VUgras{krigsfartygen}, den lilla \VUgras{undervattensbåten}
\textsuperscript{(10)} som dyker och försvinner under vattnet och också det enorma
\VUgras{pansarskeppet} \textsuperscript{(9)} som hittills fruktade nästan bara
\VUgras{torpedbåtens} \textsuperscript{(8)} anfall.

%%K t10-tit-2 sub
\VUsaut{3.0mm}
\VUpk{10.2pt}{40}{De små fartygen.}
\VUsaut{3.0mm}

%%K t10-04-1 p
4. --- Denna senare kunde redan mycket \VUgras{skickligt} finna den svaga punkten i det tjocka
metallpansaret för att där sätta in den farliga \VUgras{torpeden}, vars explosion vållar
fartygets undergång, men ändå var den mindre att frukta än undervattensbåten, för jagarna
förföljer den lätt.

%%K t10-05-1 p
5. --- Jag kunde också se de snabba pansrade \VUgras{kryssarna} \textsuperscript{(6)}, nästan
lika osårbara som det verkliga pansarskeppet, flera
\VUgras{transportfartyg} \textsuperscript{(12)} tre eller fyra
\VUgras{kanonbåtar} \textsuperscript{(7)} och slutligen en \VUgras{fregatt}
\textsuperscript{(11)}.
\VUgras{Skådespelet av den flottan, när den manövrerar för att lätta ankar är det mest intressanta.}

%%K t10-06-1 p
6. --- Vilken skillnad i byggnad mellan de stora stålmassorna och de eleganta
\VUgras{tremastarna} eller de behagfulla \VUgras{segelfartygen}
\textsuperscript{(13)} som ännu i dag används i handelsflottan och som jag såg gå in i hamnen,
för fulla segel, när jag promenerade på \VUgras{kajen} \textsuperscript{(14)}. Visserligen
förlorade dessa mycket av sina fördelar
\VUgras{när} \VUgras{vinden} var emot. De måste hala ner alla segel för att gå in
i bassängen igen och en \VUgras{bogserbåt} \textsuperscript{(16)} behövdes för att hämta dem
och föra dem, som enkla \VUgras{pråmar} \textsuperscript{(17)}, till kajen.

%%K t10-tit-3 sub
\VUsaut{3.0mm}
\VUpk{10.2pt}{40}{Handelshamnen.}
\VUsaut{3.0mm}

%%K t10-07-1 p
7. --- Det som är intressant i den hamn jag talar om, är att den inte bara innehåller en
örlogshamn, konstgjort anlagd med jättelika arbeten, utan också en underbar
\VUgras{handelshamn} som ligger i \VUgras{mynningen} av en stor flod.

%%K t10-08-1 p
8. --- Landtungan som skiljer de två bassängerna är befäst och försvarad av en rad
\VUgras{batterier} och \VUgras{bastioner} som skjuter ut i havet. Man behöver ett särskilt
tillstånd för att promenera på \VUgras{vallarna} \textsuperscript{(15)}. Jag fick det lätt,
genom min farbror, som är ingenjör vid
\VUgras{varven} \textsuperscript{(18)} och jag gick för att besöka de fartyg som
man bygger under vidsträckta skjul.

%%K t10-09-1 p
9. --- Jag bevistade till och med sjösättningen av ett pansarskepp och jag minns det gripande
ögonblick som hela folkmassan kände när den väldiga massan, överlämnad åt sig själv, började
glida på sin vaggliknande stapelbädd och kom och flöt på flodens vatten.

%%K t10-10-1 p
10. --- För att gå till varven går man över \VUgras{exercisfältet} \textsuperscript{(19)} där
garnisonens soldater varje dag kommer för att öva sig, man går över en \VUgras{liten bro}
\textsuperscript{(20)} av järn, som förbinder paradplatsen med \VUgras{arsenalen}
\textsuperscript{(21)}.

%%K t10-tit-4 sub
\VUsaut{3.0mm}
\VUpk{10.2pt}{40}{Arsenalen och dockorna.}
\VUsaut{3.0mm}

%%K t10-11-1 p
11. --- Med min farbror besökte jag den stora byggnaden full av \VUgras{kanoner}
\textsuperscript{(22)}, \VUgras{kulor} \textsuperscript{(23)} och
\VUgras{granater}. Också såg jag \VUgras{metallfabriken}
\textsuperscript{(24)} i vilken ångfartygens \VUgras{pannor} \textsuperscript{(25)}
tillverkas.

%%K t10-12-1 p
12. --- Alldeles intill är \VUgras{dockorna} \textsuperscript{(26)} --- reparationsdockor ---
och de mycket märkvärdiga \VUgras{flytdockorna} \textsuperscript{(27)} i vilka jag på mycket
nära håll, på ett fartyg som skulle repareras, kunde se en båts \VUgras{propeller}
\textsuperscript{(28)},
\VUgras{köl} \textsuperscript{(29)} och \VUgras{roder} \textsuperscript{(30)}.

%%K t10-13-1 p
13. --- Handelshamnen är lika värd att studera som örlogshamnen, och jag tillbringade många
timmar med att gapa på kajerna där \VUgras{hjulångarna} \textsuperscript{(31)} ligger förtöjda
som far uppför floden, och med att se på hur \VUgras{mastkranarna} \textsuperscript{(32)}
arbetar, som, likt fjädrar, lyfter enorma laster.

%%K t10-tit-5 sub
\VUsaut{4.0mm}
\VUpk{10.2pt}{40}{Lotsen.}
\VUsaut{3.0mm}

%%K t10-14-1 p
14. --- Jag beundrade \VUgras{atlantångarna} \textsuperscript{(34)} som gick ut från redden,
med sina \VUgras{flaggor} \textsuperscript{(35)} i full vind, förda av en skicklig
\VUgras{lots} \textsuperscript{(36)} som underbart väl vet att följa \VUgras{den lilla rännan}
som är utmärkt med \VUgras{bojar} \textsuperscript{(40)} och \VUgras{sjömärken}, medan han
undviker
\VUgras{pråmarna} \textsuperscript{(41)} som mödosamt styrs med sitt enda
fyrkantiga segel. Jag urskilde på \VUgras{förskeppet} \textsuperscript{(37)} --- fören ---
\VUgras{hytterna} \textsuperscript{(39)} i andra klass, och på
\VUgras{akterskeppet} \textsuperscript{(38)} --- aktern --- salongerna för
första klassens passagerare.

%%K t10-15-1 p
15. --- Ibland bevistade jag också lastningen av en \VUgras{pråm} \textsuperscript{(42)} i
vilken varorna från \VUgras{magasinet} \textsuperscript{(43)} och från \VUgras{sjöstationen}
\textsuperscript{(44)} just staplades, förda dit av de talrika tågen på \VUgras{kajlinjen}
\textsuperscript{(45)}.

%%K t10-tit-6 sub
\VUsaut{3.0mm}
\VUpk{10.2pt}{40}{Nöjesrodden.}
\VUsaut{3.0mm}

%%K t10-16-1 p
16. --- Ofta rodde jag i \VUgras{flytdockan} \textsuperscript{(53)}, i vilken
\VUgras{båtarna} \textsuperscript{(46)} kommer och söker skydd, och det var en
glädje för mig att hantera \VUgras{åran} \textsuperscript{(47)}. Jag klättrade upp på
\VUgras{däcket} \textsuperscript{(48)} på en \VUgras{segelbåt} \textsuperscript{(49)}, jag
klättrade, längs \VUgras{tågen} \textsuperscript{(52)}, upp i \VUgras{masten}
\textsuperscript{(51)} ända till
\VUgras{rårna} \textsuperscript{(50)}, sedan upprepade jag min tur \VUgras{i en
jolle} \textsuperscript{(54)} mellan tunga \VUgras{fiskebåtar} \textsuperscript{(55)}, där
\VUgras{näten} \textsuperscript{(56)} torkar.

%%K t10-17-1 p
17. --- På \VUgras{kajen} \textsuperscript{(58)} defilerade framför mig de mest olikartade
\VUgras{varor} \textsuperscript{(59)}, i \VUgras{lådor} \textsuperscript{(60)} eller i
\VUgras{balar} \textsuperscript{(61)}. De bildade
\VUgras{lasten} \textsuperscript{(63)} i \VUgras{skutorna}, och kraftiga
\VUgras{kranar} \textsuperscript{(62)} lyfte upp dem och satte ner dem på
\VUgras{lastvagnar} \textsuperscript{(64)} medan \VUgras{forkarlarna} anmälde dem
och betalade avgifterna på \VUgras{tullkontoret} \textsuperscript{(65)}.

%%K t10-18-1 p
18. --- Slutligen, när jag kom fram till \VUgras{svängbron} \textsuperscript{(66)} eller till
\VUgras{slussen} \textsuperscript{(69)}, medan jag gick förbi \VUgras{mudderverket}
\textsuperscript{(68)} som rensar
\VUgras{kanalen} \textsuperscript{(67)}, steg jag i land på kajen bredvid
\VUgras{semaforen} \textsuperscript{(70)}.

%%K t10-19-1 p
19. --- På andra sidan av den kajen lägger ju \VUgras{sluparna} \textsuperscript{(71)} till
som för eskaderns officerare i land. Det är ett beundransvärt skådespel att se sjömännen lyfta
sina åror i takt, medan båten som bärs av \VUgras{vågen} \textsuperscript{(72)}, lätt lägger
till vid landgångstrappan. På den platsen kan man bättre iaktta \VUgras{tidvattnet} och
rörelsen av \VUgras{flod} och \VUgras{ebb} på \VUgras{stranden} \textsuperscript{(73)}.

%%K t10-tit-7 sub
\VUsaut{3.0mm}
\VUpk{10.2pt}{40}{Sällskapshuset.}
\VUsaut{3.0mm}

%%K t10-20-1 p
20. --- För att komma hem använde jag \VUgras{den elektriska spårvagnen}
\textsuperscript{(74)} vars \VUgras{hållplats} \textsuperscript{(75)} är vid den
\VUgras{stolpe} som står framför officerarnas sällskapshus.

%%K t10-20-2 p
Från sällskapshusets \VUgras{balkong} \textsuperscript{(76)}, prydd med
\VUgras{ankare} \textsuperscript{(77)}, kanoner och kulor, såg jag ofta en
lysande samling sjöofficerare : \VUgras{löjtnanter} \textsuperscript{(78)}, med sin värja,
\VUgras{kaptener vid artilleriet} \textsuperscript{(79)} och vid
\VUgras{infanteriet} \textsuperscript{(80)} med sin \VUgras{sabel}. Bakom dem
stod en \VUgras{matros} \textsuperscript{(81)} som bar en \VUgras{kikare} till dem, när de
ville urskilja vad som sker i fjärran.

%%K t10-21-1 p
21. --- Jag såg också unga \VUgras{underlöjtnanter} \textsuperscript{(82)} som tog emot order
av sin \VUgras{kommendant} \textsuperscript{(83)} eller av
\VUgras{amiralen} \textsuperscript{(84)}, medan en \VUgras{kvartersmästare}
\textsuperscript{(85)} väntade för att föra dem till fartyget.

%%K t10-22-1 p
22. --- Från den balkongen är utsikten praktfull : man överblickar hela stranden med dess
\VUgras{tält} \textsuperscript{(86)} och \VUgras{badhytter} \textsuperscript{(87)}, och man
ser, vid badtimmen, \VUgras{badarna} \textsuperscript{(88)} som glatt leker framför
\VUgras{kasinot} \textsuperscript{(89)}.

%%K t10-23-1 p
23. --- Vid strandens ände bildar kusten en liten klippig \VUgras{halvö}
\textsuperscript{(91)}, där \VUgras{bränningen} \textsuperscript{(92)} kommer och bryter sig
och på vilken en \VUgras{fyr} \textsuperscript{(93)} är byggd, som om natten visar de
sjöfarande vägen. Den halvön är förbunden med land bara genom ett mycket smalt \VUgras{näs}
\textsuperscript{(90)}.

%%K t10-tit-8 sub
\VUsaut{3.0mm}
\VUpk{10.2pt}{40}{Allmän vy över kusterna.}
\VUsaut{3.0mm}

%%K t10-24-1 p
24. --- \VUgras{Den branta kusten} \textsuperscript{(94)} sträcker sig, uppfläkt av havet, som
nyckfullt tecknar i den en rad \VUgras{vikar} \textsuperscript{(95)} och \VUgras{uddar} (kap),
vilket gör den mycket pittoresk.

%%K t10-24-2 p
På \VUgras{platån} \textsuperscript{(96)}, som behärskar \VUgras{(havs)sundet}
\textsuperscript{(98)}, står ett mycket viktigt \VUgras{fort} \textsuperscript{(97)} och några
« villor ».

%%K t10-25-1 p
25. --- Det sundet skiljer från fastlandet en mycket farlig \VUgras{ö} \textsuperscript{(99)},
omgiven av \VUgras{skär} \textsuperscript{(100)} och
\VUgras{bränningar}, där båtar och till och med stora fartyg ibland strandar.
Trots hjälpens snabbhet, och \VUgras{livbåtarnas} snabba ankomst, är de
\VUgras{skeppsbrotten} förskräckliga, och, under dagjämningens
\VUgras{stormar}, har flera fartyg gått under med folk och varor.

%%K t10-25-2 p
Trots den grannskapen besöks \VUgras{hamnen} mycket av sjöfolket.
\VUsaut{6.0mm}
\VUfilet{22mm}
